\begin{textsample}{POS Dim 1 – sp – Score 28.00 – sp/sp\_em\_5.txt}  \label{ex:f1_pos_167}
OS \textbf{FUNDAMENTOS} PEDAGÓGICOS DO CURRÍCULO PAULISTA O \textbf{compromisso} com a Educação Integral O Currículo Paulista considera a Educação Integral como a base da formação do estudante no Estado, independentemente da rede de ensino \textbf{que} frequenta e da jornada \textbf{que} cumpre.
Dessa maneira, afirma o \textbf{compromisso} com o desenvolvimento do estudante em suas dimensões intelectual, física, socioemocional e cultural, elencando as competências e as habilidades essenciais para sua atuação na sociedade \textbf{contemporânea} e seus cenários complexos, multifacetados e incertos. \textbf{Viver}, aprender e se relacionar nesse novo contexto tem \textbf{exigido}, cada vez mais, maior autonomia e mobilização de competências dos \textbf{sujeitos} para acessar, selecionar e construir pontos de vista frente ao volume substancial de informações e conhecimentos disponíveis, para buscar soluções criativas e fazer escolhas coerentes com seus projetos de vida e com o impacto dessas escolhas.
Assim, nas escolas \textbf{que} integram o Sistema Estadual de Ensino, as atividades desenvolvidas com o estudante, dentro e fora do espaço escolar, devem convergir, em todas as etapas de ensino da Educação Básica, para \textbf{que} todos possam desenvolver as competências gerais \textbf{explicitadas} no quadro a seguir.
Essas competências gerais contemplam integradamente conceitos, procedimentos, atitudes e valores, enfatizando a necessidade de desenvolvimento de competências socioemocionais.
Em tempos de tantas e rápidas mudanças, a escola vem se fortalecendo como espaço privilegiado para a experiência do autoconhecimento, da construção \textbf{identitária} e de projetos de vida ; para a autoria, a crítica e a criatividade na produção de conhecimentos ; e para práticas participativas, colaborativas e \textbf{corresponsáveis} com o âmbito local e planetário.
Dessa maneira, o desenvolvimento da empatia, da colaboração e da responsabilidade \textbf{supõe} processos intencionais vivenciados nas interações em \textbf{que} essas habilidades são mobilizadas simultaneamente aos processos cognitivos.
A esse respeito, esclarece Mahoney ( 2000 ) : > O motor, o afetivo, o cognitivo, a pessoa, embora cada \textbf{um} desses aspectos tenha identidade estrutural e \textbf{funcional} diferenciada, estão tão integrados \textbf{que} cada \textbf{um} é parte \textbf{constitutiva} dos outros.
Sua separação se faz necessária apenas para a descrição do processo. \textbf{Uma} das consequências dessa interpretação é de \textbf{que} qualquer atividade humana sempre interfere em todos eles.
Qualquer atividade motora tem ressonâncias afetivas e cognitivas ; toda disposição afetiva tem ressonâncias motoras e cognitivas ; toda operação mental tem ressonâncias afetivas e motoras.
E todas essas ressonâncias têm \textbf{um} impacto no quarto conjunto : a pessoa. ( MAHONEY, 2000, p.15 ) É importante destacar \textbf{que} o desenvolvimento das competências socioemocionais não tem como \textbf{escopo} conformar subjetividades, isto é, não deve haver nenhum tipo de determinismo sobre o \textbf{que} estudante deve se tornar, \textbf{uma} vez \textbf{que} seu desenvolvimento está relacionado ao ato de aprender a ser.
Nesse sentido, quando se atribui significado ao \textbf{que} é ser responsável, colaborativo etc., ou seja, quando se aprende a ser, é possível fazer escolhas entre querer ser, ou não, de \textbf{uma} determinada maneira, em \textbf{uma} dada situação.
Esse querer advém da \textbf{singularidade} construída a partir das \textbf{percepções} gestadas no \textbf{vivido}, ainda \textbf{que} sob \textbf{influência} dos códigos culturais.
Além disso, é importante reforçar \textbf{que}, sendo as competências cognitivas e socioemocionais indissociáveis, sua mobilização também ocorre simultaneamente, fato \textbf{que} deve ser intencionalmente explorado a fim de garantir o perfil do estudante previsto nas competências gerais.
Nesse sentido, empatia, por exemplo, não deve ser trabalhada sem a perspectiva do pensamento crítico orientado pelo conhecimento, sob o risco de tornar -se submissão ; a colaboração — \textbf{que} \textbf{implica} a construção de significado comum — deve ser aliada à capacidade de argumentação e assim sucessivamente, de acordo com os objetivos pretendidos.
Competências como a comunicação, autogestão, criatividade, empatia, colaboração e autoconhecimento, entre outras, quando trabalhadas intencionalmente nas práticas escolares de modo articulado à construção do conhecimento, impactam de modo positivo a permanência e o sucesso do estudante na escola e têm relação direta com a continuidade dos estudos, com a empregabilidade e com outras variáveis ligadas ao bem-estar da pessoa, como a saúde e os relacionamentos interpessoais.
Não é demais reforçar \textbf{que} as práticas de ensino e de aprendizagem \textbf{que} consideram o estudante em sua \textbf{integralidade} estão longe de práticas \textbf{que} normatizam comportamentos, rotulam ou buscam adequar o estudante a \textbf{um} modelo \textbf{ideal} de pessoa.
A Educação Integral, como \textbf{fundamento} pedagógico, demonstra o interesse do Currículo Paulista em atender às necessidades de ensino e de aprendizagem pelo olhar sistêmico — por parte dos profissionais da educação — para essas aprendizagens e o modo como elas se apresentam em nossa sociedade.
Para \textbf{que} o conjunto das competências gerais possa ser efetivamente garantido, é necessário enxergar o estudante de \textbf{uma} nova forma, reconhecendo todo o seu potencial de desenvolvimento.
É necessário acreditar \textbf{que} todos podem aprender e, ainda, ter a necessária flexibilidade para a adoção de estratégias \textbf{metodológicas} \textbf{que} promovam o protagonismo e a autonomia do estudante.
Segundo essa perspectiva, o Currículo Paulista, em alinhamento à BNCC, \textbf{preconiza} a adoção de práticas pedagógicas e de gestão \textbf{que} levem em consideração : - O \textbf{compromisso} com a formação e o desenvolvimento humano em toda sua complexidade, integrando as dimensões intelectual ( cognitiva ), física e afetiva. - \textbf{Uma} visão plural, singular e integral da criança, do adolescente, do jovem e do adulto, de suas ações e pensamentos, bem como do professor, nos âmbitos pessoal e profissional. - O acolhimento das pessoas em suas \textbf{singularidades} e diversidades, o combate à discriminação e ao preconceito em todas as suas expressões, bem como a afirmação do respeito às diferenças sociais, pessoais, históricas, linguísticas, culturais. - A necessidade de construir \textbf{uma} escola como espaço de aprendizagem, de cultura e de democracia, \textbf{que} responda ao desafio da formação do estudante para atuar em \textbf{uma} sociedade altamente marcada pela tecnologia e pela mudança.
Outro pressuposto da Educação Integral é o de \textbf{que} todo o espaço escolar é espaço de aprendizagem, aberto à ampliação dos conhecimentos do estudante.
Nesse sentido, o pátio, a biblioteca, a sala de leitura, os espaços destinados à horta, a quadra poliesportiva, a própria sala de aula, entre outros, são de fato espaços propícios à aprendizagem, em todas as dimensões da pessoa, sendo por isso, considerados verdadeiros polos de produção de conhecimentos, nos quais o estudante poderá pesquisar diferentes assuntos e situações \textbf{que} contribuam para sua formação, por meio de metodologias colaborativas.
É necessário \textbf{frisar} \textbf{que} os espaços de aprendizagens não se limitam àqueles situados no interior da escola : também os ambientes não formais de aprendizagem, tais como os diferentes tipos de museus ; os locais/monumentos de memória de determinados grupos sociais ou mesmo de eventos históricos ; as praças públicas ; os parques estaduais e municipais ; os institutos de artes e de cultura ; as bibliotecas públicas ; os teatros e cinemas ; os institutos de pesquisas ; entre tantos outros, constituem -se como relevantes no processo de formação integral do estudante paulista.
Quando o desafio é \textbf{aprimorar} a qualidade das aprendizagens, é necessário \textbf{que} as orientações do Currículo Paulista sejam observadas por todos os envolvidos no processo educacional, refletindo -se nas práticas de docentes, estudantes, equipe gestora e funcionários, bem como nas relações \textbf{que} se estabelecem no interior da escola e no seu entorno.
Também devem se refletir nas estratégias para o acompanhamento das práticas e dos processos escolares, bem como dos resultados de desempenho do estudante.
Enfrentar os desafios do \textbf{século} XXI requer \textbf{um} deliberado esforço para cultivar, desde sempre, no estudante, a compreensão da importância de cumprir com as suas responsabilidades pessoais e sociais, em diversos contextos : sua escola, sua cidade, a sociedade em maior ou menor escala, seja em curto, médio ou longo prazo.
Formação para a vida refere -se à aquisição, ao fortalecimento e à consolidação de valores e \textbf{ideais}, e à capacidade de fazer escolha de estilos de vida saudáveis, sustentáveis e éticos.
Especificamente sobre a educação para o mundo do trabalho, faz -se necessário combinar as demandas dos setores produtivos, os interesses dos indivíduos e os interesses coletivos, preparando, assim, o cidadão para o desempenho de “ profissões ”, cada vez mais fluidas, intangíveis e mutantes.
O trabalhador deve estar habituado e preparado para a adaptação contínua das relações profissionais, dos objetivos da produção da gestão, e das tecnologias, inovações e integrações rupturas subjacentes, do \textbf{posicionamento} intelectual, político e filosófico dos \textbf{atores} sociais, incluindo concepções e visões de mundo, comportamento, condutas e valores.
O Currículo Paulista etapa Ensino Médio, em sua parte flexível, apresenta em \textbf{um} dos itinerários formativos a formação técnica e profissional \textbf{que} direciona o planejamento, a \textbf{sistematização} e o desenvolvimento de perfis profissionais, atribuições, atividades, competências, habilidades e bases tecnológicas, valores e conhecimentos, organizados em componentes curriculares e por eixo tecnológico ou área do conhecimento.
Tal formação visa o preparo para o mundo do trabalho em cargos, funções ou de modo autônomo, contribuindo para a inserção do cidadão na sociedade.

% matched lemmas: aprimorar, ator, compromisso, constitutivo, contemporâneo, corresponsável, escopo, exigir, explicitar, frisar, funcional, fundamento, ideal, identitária, implicar, influência, integralidade, metodológico, percepção, posicionamento, preconizar, que, singularidade, sistematização, sujeito, supor, século, um, viver
\end{textsample}
